% !TEX root = ../../../main.tex
% Sources: LEC01, IMG_9109--IMG_9111; Week 1 topic summary.
\section{Sets and countable sets}
\label{sec:analysis-i:sets-and-countable-sets}

\newthought{How do we count an infinite set?} We cannot finish listing its
elements, but we can ask whether every element has a place on a list.
Throughout these notes, $\N=\{1,2,3,\ldots\}$.

\subsection{Sets, operations, and maps}
We write $x\in A$ when $x$ belongs to $A$, and $A\subseteq B$ when every
element of $A$ belongs to $B$. Two sets are equal precisely when each is
a subset of the other. The empty set is denoted by $\varnothing$.
For sets $A$ and $B$, the basic operations are
\begin{align*}
  A\cup B&=\{x:x\in A\text{ or }x\in B\},\\
  A\cap B&=\{x:x\in A\text{ and }x\in B\},\\
  A\setminus B&=\{x:x\in A\text{ and }x\notin B\}.
\end{align*}
The complement $A^c$ is taken relative to a specified universe $U$, so
$A^c=U\setminus A$. The distributive laws and De Morgan's laws follow by
checking membership. For example,
\[
  (A\cup B)^c=A^c\cap B^c,
  \qquad (A\cap B)^c=A^c\cup B^c.
\]
For a family $(A_i)_{i\in I}$, membership in $\bigcup_{i\in I}A_i$ means
membership in at least one $A_i$; membership in $\bigcap_{i\in I}A_i$
means membership in every $A_i$.

The \emph{Cartesian product} is the set of ordered pairs
\[
  A\times B=\{(a,b):a\in A,\ b\in B\}.
\]
A map $f:A\to B$ is \emph{injective} if $f(a)=f(a')$ implies $a=a'$, and
\emph{surjective} if every $b\in B$ equals $f(a)$ for some $a\in A$.
A map with both properties is a \emph{bijection}.

\begin{definition}
  Two sets have the same \emph{cardinality}, written $|A|=|B|$, when there
  is a bijection between them. A set is \emph{countably infinite} if it
  has the same cardinality as $\N$. It is \emph{countable} if it is finite
  or countably infinite, and \emph{uncountable} otherwise.
\end{definition}
\index{countable set}\index{cardinality}
The notation $|\N|=\aleph_0$ is read ``aleph zero.'' Infinite does not
mean uncountable. For example, the even positive integers form a proper
subset of $\N$, yet $n\mapsto 2n$ is a bijection between these two sets.
\marginnote{For finite sets a proper subset has fewer elements. For infinite
sets, containment alone does not determine whether cardinality is smaller.}

\subsection{Combining countable sets}
If $A=\{a_1,a_2,\ldots\}$ and $B=\{b_1,b_2,\ldots\}$, interleave the lists
by setting
\[
  c_{2n-1}=b_n,\qquad c_{2n}=a_n.
\]
Every element of $A\cup B$ appears. If the sets overlap, discard repeated
appearances. Thus the union of two countable sets is countable, and
repeating this argument proves the same for every finite union.

What happens when there are two indices? The sum $m+n$ alone will not
distinguish $(m,n)$ from $(n,m)$. The lecture's encoding retains both pieces
of information.

\begin{proposition}
  The map $h:\N^2\to\N$ given by
  \[
    h(m,n)=m+2^{m+n}
  \]
  is injective. Consequently, the Cartesian product of two countable sets
  is countable.
\end{proposition}
\begin{proof}
  Put $d=m+n$. Since $1\leq m\leq d-1<2^d$, we have
  $2^d<h(m,n)<2^{d+1}$. These intervals are disjoint for distinct integer
  values of $d$. Hence $h(m,n)=h(j,k)$ first gives $m+n=j+k$ and then
  $m=j$, so also $n=k$. Every subset of $\N$ is countable, by listing its
  elements in increasing order. The image of $h$ therefore lists $\N^2$.
  Enumerations of two countable sets then give a list of their pairs.
\end{proof}

\begin{proposition}
  A countable union of countable sets is countable.
\end{proposition}
\begin{proof}
  Choose a list $A_m=\{a_{m,1},a_{m,2},\ldots\}$ for each nonempty set
  in the family, repeating elements if the set is finite. List the pairs
  $(m,n)$ in increasing order of $m+n$, and within each such group in
  increasing order of $m$. Every pair is reached after finitely many steps.
  Listing $a_{m,n}$ in this order reaches every element of the union.
  Discard repeated appearances to obtain the required enumeration.
\end{proof}
As usual, the argument uses a choice of enumerations for the countable
family. Figure~\ref{fig:countable-grid} shows the order of traversal.
\begin{marginfigure}[-6\baselineskip]
  \centering
  \begin{tikzpicture}[analysis diagram,x=7mm,y=7mm]
  \draw[->] (0.4,0.55) -- (4.9,0.55) node[right] {$m$};
  \draw[->] (0.55,0.4) -- (0.55,4.9) node[above] {$n$};
  \foreach \i in {1,2,3,4} {
    \node[below] at (\i,0.55) {\i};
    \node[left] at (0.55,\i) {\i};
  }
  \draw[analysis guide,->] (0.85,2.15) -- (2.15,0.85);
  \draw[analysis guide,->] (0.85,3.15) -- (3.15,0.85);
  \draw[analysis guide,->] (0.85,4.15) -- (4.15,0.85);
  \foreach \m/\n/\k in {1/1/1,1/2/2,2/1/3,1/3/4,2/2/5,3/1/6,1/4/7,2/3/8,3/2/9,4/1/10} {
    \node[fill=white,inner sep=1pt] at (\m,\n) {\k};
  }
  \node at (3.8,3.5) {$\cdots$};
\end{tikzpicture}

  \caption[Enumerating pairs along diagonals.]{Visit pairs along finite diagonals. The numerals indicate the
    order of the first ten pairs.}
  \label{fig:countable-grid}
\end{marginfigure}

\subsection{Counting the rationals}
Every positive rational has a unique reduced representation $p/q$ with
$p,q\in\N$. Thus $\Q_{>0}$ is encoded by a subset of $\N^2$ and is
countable. Adding zero and the negatives gives a countable set $\Q$.
Since $\N\subseteq\Q$, it is infinite, and therefore
\[
  |\Q|=|\N|=\aleph_0.
\]
The encoding is a map into pairs of integers; it does not assert that
rational numbers themselves are a subset of $\N$.

\subsection{Exercises on sets and enumeration}
The following exercises apply membership arguments to an arbitrary family
of sets and make the diagonal enumeration of a countable union explicit.

% !TEX root = ../../hw01.tex
% Homework 01, Exercise 1. Shared problem statement.
\begin{exercise}
  \label{ex:hw01:de-morgan}
  \solutionlink{hw01:de-morgan}
  Let $\{A_\gamma \mid \gamma \in \Gamma\}$ be a collection of sets.
  Prove De Morgan's laws
  \[
    \begin{aligned}
      \left(\bigcup_{\gamma\in\Gamma} A_\gamma\right)^c
        &= \bigcap_{\gamma\in\Gamma} A_\gamma^c,\\
      \left(\bigcap_{\gamma\in\Gamma} A_\gamma\right)^c
        &= \bigcup_{\gamma\in\Gamma} A_\gamma^c.
    \end{aligned}
  \]
\end{exercise}

% !TEX root = ../../hw01.tex
% Homework 01, Exercise 10. Shared problem statement.
\begin{exercise}
  \label{ex:hw01:countable-union}
  \solutionlink{hw01:countable-union}
  Let $A_n$ be a set such that $|A_n|=\aleph_0$ for $n=1,2,\ldots$.
  Prove $|A|=\aleph_0$ where
  \[
    A=\bigcup_{n=1}^{\infty} A_n.
  \]
\end{exercise}

